\documentclass[12pt,a4paper]{article}

\usepackage[utf8]{inputenc}
\usepackage[T1]{fontenc}
\usepackage{geometry}
\usepackage{hyperref}
\usepackage{enumitem}
\usepackage{tabularx}
\usepackage{fancyhdr}
\usepackage{titlesec}

\geometry{margin=2.2cm}
\hypersetup{
  colorlinks=true,
  linkcolor=blue,
  urlcolor=blue,
  pdftitle={IELTS Speaking Practice - User Guide},
  pdfauthor={IELTS AI},
}

\title{\Large IELTS Speaking Practice\\[2pt] \normalsize AI IELTS Speaking Assistant --- User Guide}
\author{}
\date{v0.4 --- \today}

\begin{document}

\maketitle
\thispagestyle{empty}
\vspace{8pt}
\begin{center}
  \textit{Full IELTS exam simulation | AI scoring | Live2D characters | Voice interaction}
\end{center}

\tableofcontents
\newpage

\section{Quick Start}

\subsection{Starting the Application}

\begin{itemize}[leftmargin=*]
  \item \textbf{Windows users}: Double-click \texttt{start.bat} in the project root. The script will install dependencies, start both backend and frontend, and open your browser.\\
    Visit: \texttt{http://localhost:5173}
  \item \textbf{Other systems}: See README.md for Docker or manual start instructions.
\end{itemize}

\subsection{Configuring API Key}

You need an LLM API key on first use:

\begin{enumerate}[leftmargin=*]
  \item Open the page, click \textbf{Settings} at the bottom
  \item Select your model provider from the dropdown
  \item Enter your API Key
  \item Click \textbf{Save Configuration}
\end{enumerate}

\textbf{Supported providers}:

\begin{tabularx}{\textwidth}{lX}
  \hline
  Provider & Description \\
  \hline
  DeepSeek V4 Pro/Flash & Recommended, best value \\
  OpenAI & GPT-4o-mini etc. \\
  Groq & LPU inference, very fast \\
  OpenRouter & Unified API gateway, many models \\
  Ollama & Local, no paid key needed \\
  \hline
\end{tabularx}

\subsection{First-time Tutorial}

On your first visit, a tutorial modal will appear explaining the four key steps. Click "Got it!" to dismiss. It will not appear again. You can always revisit this guide from Settings.

\newpage

\section{Exam Mode}

Exam mode simulates the full IELTS Speaking test.

\subsection{Exam Flow}

\begin{enumerate}[leftmargin=*]
  \item \textbf{AI examiner introduction} --- An AI examiner introduces the test rules.
  \item \textbf{Part 1: Introduction \& Interview (4--5 min)} --- The examiner asks 3 questions around one topic (randomly chosen from 43 topics including Hometown, Work/Study, Technology, etc.).
  \item \textbf{Part 2: Individual Long Turn} --- A Cue Card appears with a topic and 4 prompt points. You have 60 seconds to prepare, then speak for 1--2 minutes. The examiner may ask 1--2 follow-up questions.
  \item \textbf{Part 3: Two-way Discussion} --- The examiner asks 2--3 deeper discussion questions related to the Part 2 topic.
  \item \textbf{Exam ends} --- The examiner announces the end; the system navigates to the scoring report automatically.
\end{enumerate}

\subsection{Controls}

\begin{itemize}[leftmargin=*]
  \item Click the \textbf{Microphone button} at the bottom to start recording
  \item Recording auto-stops when the timer expires; you can also stop manually
  \item A Live2D virtual examiner character (Haru) appears on screen, with expressions and gestures
\end{itemize}

\subsection{Timer Reference}

\begin{tabularx}{\textwidth}{lX}
  \hline
  Stage & Default Time \\
  \hline
  Part 1 each question & 45 seconds \\
  Part 2 preparation & 60 seconds \\
  Part 2 speaking & 120 seconds \\
  Part 2 follow-up & 60 seconds \\
  Part 3 each question & 60 seconds \\
  \hline
\end{tabularx}

\newpage

\section{Scoring Report}

After the exam, a full scoring report is generated:

\subsection{Scoring Criteria}

Based on the official IELTS four criteria (0--9, 0.5 increments):

\begin{itemize}[leftmargin=*]
  \item \textbf{Fluency and Coherence} --- Speech rate, pauses, logical flow
  \item \textbf{Lexical Resource} --- Vocabulary range, precision, paraphrasing
  \item \textbf{Grammatical Range and Accuracy} --- Sentence complexity, error frequency
  \item \textbf{Pronunciation} --- Clarity, stress, intonation (estimated from text quality)
\end{itemize}

\subsection{Overall Band Score}

Average of the four criteria, rounded to the nearest 0.5.

\subsection{Suggestions}

The report includes 3 personalised suggestions based on your performance.

\vspace{8pt}
\begin{center}
  \fbox{\parbox{0.9\textwidth}{\small
    \textbf{Note}: Pronunciation scoring is an indirect estimate from the transcribed text. For more accurate pronunciation feedback, enable Whisper local ASR in Settings.
  }}
\end{center}

\newpage

\section{Free Chat Mode}

Free Chat mode has no scoring and no time limits. It is designed for casual English practice.

\begin{itemize}[leftmargin=*]
  \item The AI greets you in English and begins a conversation
  \item Click the microphone to speak; the AI replies in voice
  \item You can change topics freely
  \item Chat history is not saved (cleared on refresh)
  \item Live2D character is Mao (casual style), different from Haru (formal, exam mode)
\end{itemize}

\section{Voice Input}

\subsection{Browser ASR (Default)}

Uses Chrome/Edge built-in Web Speech API. No extra configuration needed. Low latency, works out of the box. However, accuracy depends on the browser, and continuous recording may stop after $\sim$60 seconds.

\subsection{Whisper Local ASR}

Uses faster-whisper for local speech recognition. Higher accuracy and supports longer recordings.

\textbf{How to enable}:
\begin{enumerate}[leftmargin=*]
  \item Go to Settings $\to$ Whisper ASR
  \item Toggle \textbf{Enabled}
  \item Download a model (recommended: \texttt{small}, $\sim$460\,MB)
  \item Select your language
  \item Save configuration
\end{enumerate}

Whisper runs on CPU. Supported models range from tiny (75\,MB) to large-v3 (2.9\,GB) -- 9 models total.

\section{Appearance and Language}

\subsection{Live2D Characters}

Two built-in Live2D virtual characters:

\begin{itemize}[leftmargin=*]
  \item \textbf{Haru (Exam mode)} --- Formal style, nods, smiles, gestures
  \item \textbf{Mao (Free Chat mode)} --- Casual, expressive
\end{itemize}

\subsection{Eye Tracking}

In Settings, you can set the character's gaze:
\begin{itemize}[leftmargin=*]
  \item \textbf{Follow Mouse} --- The character's eyes follow your mouse
  \item \textbf{Look Forward} --- The character looks straight ahead
\end{itemize}

\subsection{UI Language and Training Language}

\begin{itemize}[leftmargin=*]
  \item \textbf{UI Language}: Switch display language (English / Chinese)
  \item \textbf{Training Language}: Controls speech recognition target language and TTS voice. Supports English, Chinese, Japanese, Korean, French, German, Spanish (7 languages)
\end{itemize}

\newpage

\section{Troubleshooting}

\subsection{AI doesn't respond?}

Check:
\begin{enumerate}[leftmargin=*]
  \item API Key is correctly entered in Settings
  \item API account has sufficient credits
  \item Network connection is working (some services may require a VPN)
\end{enumerate}

\subsection{Voice recognition not working?}

\begin{enumerate}[leftmargin=*]
  \item Ensure your browser has microphone permission
  \item Use Chrome or Edge (Firefox/Safari do not support Web Speech API)
  \item If browser ASR frequently disconnects, try enabling Whisper local ASR
\end{enumerate}

\subsection{Port already in use?}

The application uses port 8000 (backend) and 5173 (frontend). Close any programs using these ports and restart.

\subsection{Browser didn't open automatically?}

Manually navigate to \texttt{http://localhost:5173} in your browser.

\subsection{How to reset the tutorial?}

Open your browser's developer console (F12) and run:
\begin{verbatim}
localStorage.removeItem("tutorialSeen")
\end{verbatim}
Then refresh the page.

\subsection{How to add new exam types or languages?}

See the ``Adding a new exam'' and ``Adding a new language'' sections in README.md.

\section{Technical Information}

\begin{tabularx}{\textwidth}{lX}
  \hline
  Component & Technology \\
  \hline
  Frontend & React 19 / Vite 8 / TypeScript / PixiJS 6 \\
  Backend & Python FastAPI / Uvicorn \\
  LLM & OpenAI SDK compatible (DeepSeek / OpenAI / Groq / OpenRouter / Ollama) \\
  ASR & Web Speech API / faster-whisper \\
  TTS & Web Speech API (Edge natural voices) \\
  Live2D & pixi-live2d-display + Cubism 4 SDK \\
  \hline
\end{tabularx}

\vspace{1.5cm}
\begin{center}
  \textit{--- End of User Guide ---}\\[4pt]
  GitHub: \url{https://github.com/Making-acid/ielts-speaking-ai}
\end{center}

\end{document}
